% BHU PYQ -- Manually transcribed & error-corrected
% Paper: GLB-602: Stratigraphy
% Exam: B.Sc. (Hons.) Semester VI Examination, 2022-23

\documentclass[11pt,a4paper]{article}
\usepackage[a4paper,top=2.5cm,bottom=2.5cm,left=2.8cm,right=2.8cm,headheight=14pt]{geometry}
\usepackage{amsmath,amssymb}
\usepackage[T1]{fontenc}
\usepackage[utf8]{inputenc}
\usepackage{lmodern,microtype,fancyhdr,enumitem,hyperref,lastpage,parskip}

\hypersetup{colorlinks=true,urlcolor=black,linkcolor=black,
  pdftitle={GLB-602 Stratigraphy -- BHU B.Sc. Sem VI 2022-23}}
\pagestyle{fancy}\fancyhf{}
\renewcommand{\headrulewidth}{0.4pt}\renewcommand{\footrulewidth}{0.4pt}
\fancyhead[L]{\small   \textit{Banaras Hindu University}}
\fancyhead[R]{\small   \textit{GLB-602: Stratigraphy}}
\fancyfoot[C]{\small Page  \thepage\ of \pageref{LastPage}}


\newlist{parts}{enumerate}{2}
\setlist[parts,1]{label=(\alph*),leftmargin=*,itemsep=5pt,topsep=3pt}

\newcommand{\pts}[1]{\hfill   \textit{[#1]}}


\begin{document}

\begin{center}
  {\large\bfseries BANARAS HINDU UNIVERSITY}\\[2pt]
  {
\normalsize B.Sc. (Hons.) --- Semester VI Examination, 2022-23}\\[6pt]
  \rule{0.8\linewidth}{0.5pt}\\[6pt]
  {\large\bfseries Paper: GLB-602 --- Stratigraphy}\\[4pt]
  {
\normalsize Time: 3 Hours \hfill Maximum Marks: 70}\\[2pt]
  {\small Ref. No.: 058441}
\end{center}
\rule{\linewidth}{0.4pt}
\smallskip



\noindent   \textit{Instructions: Answer   \textbf{five} questions in all, selecting at least   \textbf{two} each from Sections A and B, and Section-C which is compulsory. All questions are of equal value.}
\bigskip

\bigskip


\noindent {\large\bfseries SECTION --- A}
\rule{\linewidth}{0.4pt}\smallskip




\noindent   \textbf{Question 1.} Discuss ``principles of stratigraphy'' and various criteria for stratigraphic classification and correlation. \pts{14}

\medskip


\noindent   \textbf{Question 2.} Write short notes on   \textbf{any two} of the following: \pts{$2  \times 7 = 14$}

\begin{parts}
  \item Neogene-Quaternary stratigraphic successions of the Siwalik Supergroup
  \item Sequence stratigraphy and concepts of accommodation space
  \item Lower Triassic of Spiti
\end{parts}

\medskip


\noindent   \textbf{Question 3.} Discuss the stratigraphic framework of the Jurassic of Kachchh. \pts{14}

\medskip


\noindent   \textbf{Question 4.} Give a detailed account of the Paleogene stratigraphic successions of Himachal Pradesh and Assam. \pts{14}

\bigskip


\noindent {\large\bfseries SECTION --- B}
\rule{\linewidth}{0.4pt}\smallskip




\noindent   \textbf{Question 5.} Discuss in detail the Paleozoic stratigraphic successions of Kashmir. \pts{14}

\medskip


\noindent   \textbf{Question 6.} Give brief accounts of   \textbf{any two} of the following: \pts{$2  \times 7 = 14$}

\begin{parts}
  \item Origin and stratigraphic age of the Deccan Traps
  \item Stratigraphic successions of Upper Gondwana
  \item Karewa Group
\end{parts}

\medskip


\noindent   \textbf{Question 7.} Write short notes on   \textbf{any four} of the following: \pts{$4  \times 3\frac{1}{2} = 14$}

\begin{parts}
  \item Neogene/Quaternary boundary
  \item Zewan Formation
  \item Bhuj Formation
  \item Anaipadi Formation
  \item Baratang Formation
  \item Stratotype and type locality
\end{parts}

\bigskip


\noindent {\large\bfseries SECTION --- C}
\rule{\linewidth}{0.4pt}\smallskip




\noindent   \textbf{Question 8.} Fill in the blanks: \pts{$7  \times 2 = 14$}

\begin{parts}
  \item The stratigraphic age of the Syringothyris Limestone is \underline{\hspace{4cm}}.
  \item The Cretaceous/Paleocene boundary lies at \underline{\hspace{3cm}} My.
  \item The age of the Umaria marine bed is considered as \underline{\hspace{4cm}}.
  \item The stratigraphic age of the Productus Limestone is considered as \underline{\hspace{4cm}}.
  \item The Triassic/Jurassic boundary lies at \underline{\hspace{3cm}} My.
  \item ``Stage'' is a \underline{\hspace{4cm}} unit.
  \item GSSP stands for \underline{\hspace{4cm}}.
\end{parts}

\vfill\rule{\linewidth}{0.4pt}

\begin{center}\small
\href{https://bhu-exam.vercel.app/}{Click here to visit our website}\\[4pt]
\href{https://me-aryan.vercel.app/}{Click here to visit our contributor}\\[4pt]
\href{https://quashan-me.vercel.app/}{Click here to visit our contributor}
\end{center}
\end{document}
